% Chapter 6: Starting New Projects with Claude
\chapter{Starting New Projects with Claude}
\label{ch:greenfield}

Starting a new project with Claude is the easiest scenario to get right
and the most common one to get wrong. Easy because there is no legacy
context to manage. Wrong because developers skip configuration and treat
Claude like a chatbot instead of a development environment.

This chapter provides a day-by-day protocol for the first month.

\section{Greenfield: From Zero to Production}
\label{sec:greenfield-protocol}

\marginnote[Official]{Run \code{/init} to auto-generate CLAUDE.md. Claude analyzes your repo and proposes conventions.\sourceurl{https://code.claude.com/docs/en/best-practices}}

\subsection{Day 1: Foundation}

\begin{enumerate}
  \item \textbf{Initialize the project} with your language's tooling
    (\code{npm init}, \code{cargo init}, \code{poetry init}).
  \item \textbf{Run \code{/init}} to generate a starter CLAUDE.md.
    Claude examines your project structure and proposes build commands,
    style rules, and architecture notes.
  \item \textbf{Set up permissions} in \code{.claude/settings.json}.
    At minimum: deny access to \code{.env*} and \code{secrets/}.
  \item \textbf{Add one hook}: \code{PostFileWrite $\rightarrow$ auto-format}.
    This prevents style drift from the first file onward.
\end{enumerate}

\convergence{Start with minimal CLAUDE.md (build commands + code style),
add rules/skills/hooks as patterns stabilize.
Do not over-engineer configuration before you know what you need.}

\subsection{Week 1: Patterns Emerge}

\begin{enumerate}
  \item \textbf{Add testing infrastructure.}
    Configure test runner in CLAUDE.md: ``Run \code{pytest tests/} after changes.''
  \item \textbf{Add your first custom command} for the most common workflow.
    Typically \code{/test} or \code{/check}.
  \item \textbf{Add conditional rules} if the project has distinct modules.
  \item \textbf{Configure PreCommit hook} $\rightarrow$ run tests.
    This is the single most valuable hook for new projects.
\end{enumerate}

\subsection{Month 1: Stabilization}

\begin{enumerate}
  \item \textbf{Review and trim CLAUDE.md.}
    Remove instructions Claude consistently follows without being told.
    Add instructions for rules Claude keeps violating.
  \item \textbf{Add skills} for repeated multi-step workflows.
  \item \textbf{Set up memory} for cross-session persistence of key decisions.
  \item \textbf{Document architecture} in CLAUDE.md---not as prose,
    but as actionable instructions: ``API routes go in \code{src/routes/}.
    Each route has a corresponding test in \code{tests/routes/}.''
\end{enumerate}

\margintip{Most solo projects stabilize at Tier 2--3 configuration within one month. Resist the urge to jump to Tier 5.}

% -------------------------------------------------------------------
\section{Quality Gates from Day One}
\label{sec:quality-gates}

\convergence{Verification criteria in CLAUDE.md from the start.}
This is the single configuration that most improves project outcomes.

\begin{minted}{markdown}
## Quality Standards
- Always run tests after code changes
- Never commit without passing lint + tests
- Include edge case tests for every new function
- Coverage target: 80% for src/ modules
\end{minted}

\subsection{Hook-Enforced Standards}

\practitioner{Three hooks that prevent quality debt before it accumulates:}

\begin{minted}{json}
{
  "hooks": {
    "PostFileWrite": [{
      "command": "prettier --write \"$FILE\"",
      "description": "Auto-format every written file"
    }],
    "PreCommit": [{
      "command": "npm run lint && npm test",
      "description": "Block commits with lint errors or test failures"
    }],
    "PostCommit": [{
      "command": "npx coverage-diff",
      "description": "Report coverage change per commit"
    }]
  }
}
\end{minted}

\subsection{Phase-Appropriate Rigor}

Recall from \cref{sec:phases}: the exploration phase is acceptable for the
first sprint. But set explicit graduation criteria in CLAUDE.md:

\begin{minted}{markdown}
## Phase: EXPLORATION (until March 15)
After prototype validated:
- Switch to DEVELOPMENT phase
- Require 80% coverage, type hints
- Enable strict linting
\end{minted}

% -------------------------------------------------------------------
\section{When Claude Writes the First Draft}
\label{sec:test-first}

\marginnote[Official]{``Give Claude a way to verify its own work.'' For greenfield, this means: provide test cases in your initial prompt, not after the code is written.\sourceurl{https://code.claude.com/docs/en/best-practices}}

\practitioner{The test-first pattern produces higher-quality greenfield code
than code-first:}

\begin{enumerate}
  \item Describe the \textbf{interface}: inputs, outputs, edge cases.
  \item Claude writes \textbf{tests} based on the interface.
  \item Claude writes \textbf{implementation} that passes the tests.
  \item Tests run \textbf{automatically} via hooks.
\end{enumerate}

\begin{minted}{text}
Prompt: "Create a rate limiter class with these requirements:
- Token bucket algorithm
- Configurable rate (requests/second) and burst size
- Thread-safe
- Returns (allowed: bool, retry_after: float)

Write tests first, then implementation."
\end{minted}

This works because Claude excels at test generation when given
clear specifications. The tests then constrain the implementation,
preventing the ``looks correct but is subtly wrong'' failure mode
described in \cref{ch:testing}.

\begin{keyinsight}
The first week of a new project determines its trajectory.
Three investments on day one---\code{/init}, a deny rule for secrets,
and a PreCommit test hook---compound into dramatically better outcomes
over the project's lifetime.
\end{keyinsight}
